\section{Set the interfaces before allocating registers}
The project uses little-endian MARS MIPS64 Release 6. Addresses and stack saves are
64-bit; algorithm words are 32-bit. All functions return through the supplied
\texttt{ra}, preserve every bit of \texttt{s0} through \texttt{s7}, \texttt{fp} and
\texttt{sp}, and return a status in \texttt{v0}. The stack is eight-byte aligned.
Temporary scalar registers may change. The MSA block function may also change
\texttt{w0} through \texttt{w9}, including the floating-point storage they alias.
Do not retain a live floating-point value in that storage across the call.

\begin{description}
  \item[\texttt{chacha\_quarter(a0,a1,a2,a3)}]
        Each argument points to a distinct, word-aligned writable word, in the order
        \(a,b,c,d\). Replace those four words with one quarter-round's result. Their addresses
        need not be adjacent. Return zero; access no other caller bytes.
  \item[\texttt{chacha\_block\_scalar(a0,a1)}]
        Read sixteen word-aligned input words at \texttt{a0}; preserve them. Write exactly
        64 output bytes at the disjoint, word-aligned \texttt{a1}. The input is an already
        constructed state. Apply ten double-rounds, add the original input words, and serialize
        the sums in little-endian order. Return zero.
  \item[\texttt{chacha\_block\_msa(a0,a1)}]
        Use the same input, output and preservation contract as the scalar block, with MSA
        word lanes. Four-byte alignment is sufficient for this project's \texttt{ld.w} and
        \texttt{st.w} accesses; a 16-byte vector does not make every vector access require
        16-byte alignment.
  \item[\texttt{chacha\_xor(a0,a1,a2,a3)}]
        The unchanged state is at \texttt{a0}, the input bytes at \texttt{a1}, output bytes at
        \texttt{a2}, and the unsigned byte count in \texttt{a3}. This local interface accepts
        0 through 4096 bytes. The state is word-aligned; byte buffers may be unaligned. Input
        and output may be identical or disjoint. Other overlap is outside the interface.
        Return zero on success, one for an excessive count, or two if the required counters
        would exceed \texttt{ffffffff}. Reject before writing output. Check the count first;
        a zero count returns without accessing any supplied pointer. For a positive valid
        count the state and buffers must be mapped, and the state is disjoint from both buffers.
\end{description}

The bounded XOR interface makes its whole-request counter check easy to inspect.
4096 is a teaching limit, not an algorithm limit. An error leaves caller memory
unchanged. A successful call does not update the input counter. A second call with
that same state restarts at the same keystream position; the caller must deliberately
select the correct next counter or keep a richer streaming context. There is no
hidden global state that advances on behalf of the caller.
